\chapter{Short-range forces: P$^3$M and TreePM}\label{ch:short-range}
\chapterstatus{written}{2026-11-14}

A PM force is accurate only on scales several cells above the mesh spacing. To resolve gravity below the mesh scale,
\fml{} splits the force into a long-range part, computed on the mesh, and a short-range part, summed over particle
pairs within a cutoff: the P$^3$M method with cell lists, or a TreePM method with a monopole octree. The same split
supports Newtonian gravity plus a Yukawa fifth force, the linear force of chameleon-like theories. This chapter
describes the kernels, the search structures, the MPI ghost exchange and the integration into SimulationCode.

\begin{modulebox}{ShortRangeGravity}
Namespace \nsp{FML::SHORT_RANGE_GRAVITY}, header-only: \file{ShortRangeGravity.h} (\cls{Options}, \cls{Kernel},
\code{direct_sum}, \cls{CellList}), \file{Softening.h}, \file{Yukawa.h}, \file{Octree.h} (\cls{Octree}),
\file{LongRangePM.h} (\code{apply_long_range_filter}), \file{SlabGhosts.h} (\code{exchange_slab_ghosts}).\par
SimulationCode: \file{src/ShortRangeForce.h}, \file{src/ShortRangeInteraction.h}, \file{src/ShortRangeSettings.h}.\par
Documentation: \file{ShortRangeGravity/README.md}, \file{ShortRangeGravity/tests/RESULTS.md},
\file{SimulationCode/SHORT_RANGE.md}, \file{COLA_SHORT_RANGE.md}, \file{SHORT_RANGE_RESULTS.md},
\file{CLUSTERED_FORCE.md}, \file{DYNAMICS.md}.
\end{modulebox}

\section{The force split}\label{sec:sr-split}

With a Gaussian long-range filter of scale $r_s$ (the convention of GADGET-2 \cite{Springel2005}),
\begin{equation}
\phi_{\rm long}(\vk)=\phi(\vk)\,e^{-k^2r_s^2},
\end{equation}
the complementary pair acceleration of a point mass is
\begin{equation}
\bm a_{\rm short}=Gm\frac{\vx_j-\vx_i}{r^3}\left[\mathrm{erfc}(q)+\frac{2q}{\sqrt\pi}e^{-q^2}\right],\qquad q=\frac{r}{2r_s},
\end{equation}
which decays like a Gaussian beyond a few $r_s$. \code{apply_long_range_filter} multiplies a dedicated Fourier copy of
the density by the filter (in box units, $r_s/B$) before the ordinary PM solve. The analytic split does not remove the
mesh's anisotropy and aliasing: the long force must be resolved, i.e.\ $r_s$ of order one to two mesh cells or more.

\paragraph*{Softening.} Particles have the compact cubic-spline density of support $h$
(\param{force_short_range_softening_support}),
\begin{equation}
W(r)=\frac{8}{\pi h^3}\begin{cases}1-6u^2+6u^3, & u<\frac12,\\ 2(1-u)^3, & \frac12\le u<1,\end{cases}\qquad u=r/h,
\end{equation}
so the force is exactly Newtonian beyond $h$. $h$ is the support radius, not the Plummer-equivalent softening:
$h=2.8\,\epsilon$. Inside $h$ the Gaussian long-range force is subtracted from the full spline-softened force
(multiplying the softened force by the unsoftened short window would give a different total force); a small-$r$
series avoids cancellation, and coincident particles give zero force.

\paragraph*{Cutoff.} \param{force_short_range_cutoff_tolerance} sets the radius where the short-range pair force,
relative to unsoftened Newtonian gravity, falls below the tolerance; the force is truncated there without tapering.
This bounds a single pair's window, not the relative error of a net force near a cancellation. One requires
$h\le r_{\rm cut}<B/2$. The pair coefficient is tabulated on $[h,r_{\rm cut}]$ with linear interpolation
(\param{force_short_range_table_points} $=16384$; error $<10^{-8}$ of $G/r^3$, $3\times$ faster on clustered
catalogues; 0 evaluates it exactly).

\section{Newtonian plus Yukawa gravity}\label{sec:sr-yukawa}

A light scalar of mass $m=1/\lambda$ coupled to matter adds a Yukawa term to the point-mass potential,
\begin{equation}
\Phi=-\frac{Gm}{r}\left[1+\alpha e^{-r/\lambda}\right],\qquad
\text{force}\propto1+\alpha\left(1+\frac r\lambda\right)e^{-r/\lambda},
\end{equation}
(\code{Kernel(options, NewtonYukawa{alpha, lambda})}; $\alpha$ of either sign). The long-range Fourier kernel is
\begin{equation}
\frac{e^{-k^2r_s^2}}{k^2}+\alpha\frac{e^{-(k^2+\mu^2)r_s^2}}{k^2+\mu^2},\qquad \mu=1/\lambda,
\end{equation}
the Helmholtz-aware Ewald split of the Yukawa potential (Eqs.~30--32 of the reference in \file{REFERENCES.md},
doi:10.1103/PhysRevE.103.043204), and the complementary short-range Green's function is
\begin{equation}
\frac{e^{-\mu r}\,\mathrm{erfc}(q-z)+e^{\mu r}\,\mathrm{erfc}(q+z)}{2r},\qquad q=\frac{r}{2r_s},\quad z=\mu r_s,
\end{equation}
evaluated with scaled erfc functions, series and heat-kernel integrals for the extreme limits, and stable differences
for negative $\alpha$. Yukawa softening is a pair-potential prescription: with $s(r)=1/g_N(r)$ from the softened
Newtonian Green's function, the softened Yukawa term is $e^{-s/\lambda}/s$, which is finite at the centre and exact
outside the support. The cutoff conservatively solves $(1+|\alpha|)\times$ Newtonian window $=$ tolerance. All lengths
are in the coordinate unit: for a fixed physical range in comoving coordinates pass $\lambda_{\rm com}=\lambda_{\rm phys}/a$.

\section{Search structures and parallelism}\label{sec:sr-search}

\begin{itemize}
\item \code{direct_sum}: the reference $\mathcal O(N^2)$ sum.
\item \cls{CellList} (P$^3$M): contiguous cells at least as wide as the cutoff; \code{max_cells} bounds the metadata
      (coarser cells, more pair work, never missed neighbours). Pairwise momentum is conserved to summation precision.
\item \cls{Octree} (TreePM): flat nodes with mass, centre of mass and bounds, monopole acceptance. A node is accepted if
      it is entirely inside the cutoff, outside the softening, within one periodic image, does not contain the target,
      and is small relative to both the distance and the split scale (\param{force_short_range_opening_angle}, $0.3$ a
      starting value; $0$ gives exact leaf sums). With Yukawa forces the opening also resolves $\lambda$. One-sided
      tree approximations need not conserve momentum exactly.
\item OpenMP over targets with no shared writes; results are attractive accelerations with source mass and coupling
      included; targets can be a subset of the sources; self-interactions are excluded by ID.
\item \code{exchange_slab_ghosts}: each source is sent once to every remote slab within the cutoff (including periodic
      and non-adjacent slabs), in bounded chunks with 64-bit counts; no all-gather of particles. Ghosts and search
      structures are rebuilt after every drift.
\end{itemize}

\paragraph*{Units in \fml{}.} With box coordinates $x=r/L$, $\nabla_x^2\Phi=C\delta$ and total source mass $M$, the pair
potential is $\Phi_{\rm pair}=-\frac{C}{4\pi}\frac mM\frac{L}{|r-r_j|}$, so the coupling of the short force is
\code{fml_acceleration_coupling(C, M, L)} $=CL^2/(4\pi M)$, with $C=\frac32\Om_ma$ in SimulationCode's GR and $M$ the
global source mass.

\begin{cppcode}[P$^3$M and TreePM short-range forces]
#include <FML/ShortRangeGravity/ShortRangeGravity.h>
#include <FML/ShortRangeGravity/Octree.h>
#include <FML/ShortRangeGravity/SlabGhosts.h>
using namespace FML::SHORT_RANGE_GRAVITY;
Options options;
options.boxsize = 250; options.split_scale = 2; options.softening_support = 0.1;
options.cutoff_tolerance = 1e-6; options.coupling = 1;
Kernel kernel(options);                                  // or Kernel(options, NewtonYukawa{1.0/3, 10})
auto sources = exchange_slab_ghosts(owned, xmin, xmax, kernel);
auto p3m     = CellList(sources.particles, kernel).compute(owned);
auto treepm  = Octree(sources.particles, kernel, TreeOptions{0.3, 16, 40}).compute(owned);
\end{cppcode}

\section{Short-range forces in SimulationCode}\label{sec:sr-simcode}

With \param{force_short_range} \code{= true} the force of a step is the filtered PM force plus the short-range kick,
at a common time step:
\begin{equation}
\bm v_{\rm new}=s\,\bm v_{\rm old}-\Delta_K\,\bm F_{\rm long}+\Delta_K\,\bm A_{\rm short},
\end{equation}
(the mesh returns $+\nabla\phi$, the pair code an acceleration), with the short kernel normalised by $C(a)/(4\pi M)$ times
the model coupling. With force interlacing only the long force is averaged over the shifted copies; one unshifted
short kick follows. The lightcone potential is computed from the unfiltered density. In COLA the residual velocity is
updated as $\bm u'=s\bm u+\Delta_K(\bm A_{\rm long}+\bm A_{\rm short}-\bm A_{\rm LPT})$, with the frame potentials
unfiltered.

\paragraph*{Modified gravity.} A gravity model describes its pair potential outside the softening as
\begin{equation}
\Phi_{\rm pair}=-\frac{GM}{r}\left[\mu_0+(\mu_\infty-\mu_0)e^{-r/\lambda}\right],\qquad
\mu_0=G_{\rm eff}(a,0)/G,\quad \mu_\infty=G_{\rm eff}(a,k\to\infty)/G,
\end{equation}
with the comoving range \code{short_range_comoving_range(a)} ($0$: infinite; $f(R)$, the symmetron and $(m,\beta)$ give
their inverse mass divided by $a$). The mesh keeps the remainder
$\mu_{\rm mesh}(k)=G_{\rm eff}(a,k)/G-\mu_{\rm pair}(k)+\mu_{\rm pair,long}(k)$, which vanishes for the analytic Yukawa
models. Nonlinear screening has no density-independent pair kernel, so screened and exact solvers are rejected with
short-range forces.

\paragraph*{Scope.} Short-range forces support 3D unscreened gravity with one cold species, ordinary $N$-body with Quinn
steps or 1LPT/2LPT COLA with Quinn or Tassev steps, the continuous Green's function and gradient, and assignment-based
densities. Not yet: subcycling, massive neutrinos or separate species, the $N$-body gauge, FastPM/BullFrog, NUFFT/sheet
or early-time forces, and an acceleration-based time-step controller (short-range forces can need many more steps).

\begin{paramfile}[SimulationCode: PM + short-range forces]
force_short_range                   = true
force_short_range_backend           = "cell"   -- or "tree"
force_short_range_split_scale       = 3.125    -- comoving Mpc/h: two cells of a 64^3 mesh in 100 Mpc/h
force_short_range_softening_support = 0.2      -- spline support h = 2.8 epsilon, comoving Mpc/h
force_short_range_cutoff_tolerance  = 1e-6
force_greens_function_kernel        = "continuous"
force_gradient_kernel               = "continuous"
timestep_method                     = "Quinn"
\end{paramfile}

\begin{validation}
\begin{itemize}
\item Kernels: independent spline enclosed-mass quadrature; split and zero-separation limits; filtered PM plus short
      force against independent periodic Ewald sums (worst relative error $\sim3\times10^{-6}$); Yukawa kernels against
      100-digit Ewald references and image sums ($1.7\times10^{-13}$); unit invariance.
\item Search: cell lists and exact octrees equal direct sums to summation precision (uniform, clustered, periodic,
      coincident, zero-mass); ghost completeness for wide cutoffs, empty and non-adjacent slabs.
\item SimulationCode: the actual PCS deposition, force and short correction against Ewald sums for an off-grid
      clustered fixture: relative RMS force error 0.59 for unsplit PM on $16^3$, and $1.4\times10^{-3}$,
      $6\times10^{-5}$, $2\times10^{-6}$ for the combined force on $16^3$, $32^3$, $64^3$ meshes; second-order time-step
      convergence of a 512-particle run; tree and cell trajectories identical; identical states on 1, 2 and 4 ranks.
\item Dynamics: two-body orbits and 16-particle clusters through the real simulation step against adaptive direct
      Ewald evolution: errors $\sim10^{-3}$ of the initial radius and velocity dispersion with 256--1024 steps,
      converging at second order (\file{DYNAMICS.md}); clustered static-force error distributions
      (\file{CLUSTERED_FORCE.md}).
\end{itemize}
These are controlled verification cases, not a calibration of nonlinear cosmological statistics.
\end{validation}
